
Challenge 10B. If you’re new, you’ll soon find out. If you’re one of a returning participant, you know the deal. 
One of the most difficult things I find with Challenge 10B every year, apart from the cipher itself, is taking a break. In the midst of chaotic, spiralling ideas on how to crack the final cipher, it’s so easy to fall into that trap of “I’m so close… just a bit longer…” and then 4 hrs later your eyes hurt, brain’s fried, and you’re resisting the urge to throw your laptop out the window. 
It feels completely counterintuitive, but stepping away really does help. And even though I know that, I still find it hard to move away when I’m in the zone. I just wanted to write this post because I’m sure a lot of you out there can also relate to this feeling, so I thought I’d share a few things that have helped me reset and come back with an insight for 10B, and maybe they might help you too: 
– Write down what you know – every hour or just pick a timeframe, jot down everything you’ve figured out about the cipher, even the obvious stuff, especially the obvious stuff. Key lengths/repeated n-grams/ letter positions/cipher text arrangement/cribs, etc. It’s easy to feel like you’re not making much progress, but when you see it on paper, it’s surprising how much you’ve actually uncovered, and it helps track progress
– Have a snack/remember to eat – the brain burns a ridiculous amount of energy when you’re problem-solving, and as I’ve been told many times in training: food is fuel.
– Check in with your team if you have one – it’s easy to disappear into your own rabbit hole of ideas so it’s always useful take a moment to sync and share weird hunches, and remind each other to take breaks as well.
– Setting expectations at home if need be – my family is always made well aware that I will be busy during cipher challenge season, especially at the end so know what to/not to expect from me. I find that this definitely does reduce some background stress 
– Slightly unconventional, but I have a special hoodie I only wear when working on the cipher challenge (sadly it doesn’t have any actual powers, but having a dedicated codebreaking hoodie does make me feel extra geeky). When I take a break, I physically change out of it. It sounds silly but that very act of changing the hoodie helps me switch mental gears and actually rest. As much as we all love the challenge, I’m sure we all have other things (including food and sleep) that need getting done and I just found that having a physical reminder helped me compartmentalise 
This is my fifth year doing the challenge, and honestly, my only real regret was not sticking with 10B till the end in my second year. I gave up too early and convinced I was going nowhere after working on it for just over a day, only to see others break through with just a bit more persistence. What I’ve learned since then is that hints do come, but you’ve got to have patience. They’re mostly be subtle or buried under obscure forum comments, or phrased in a way that only makes sense after you’ve cracked it. 
It’s funny but I think that if you’re annoyed at this cipher, that’s kind of amazing in itself. It means you’ve found something that you care about, something that’s gripped your attention so much you’re still here, thinking and trying, and that’s rare. Most people go through life bored by what they do because let’s be real, if a bunch of weirdly spaced letters and numbers are making you want to throw your keyboard across the room, maybe that’s just be a sign you’ve found something worth fighting with (and for), and that’s a great thing! 
Well done to anyone who’s cracked 10B, and equally well done to anyone who’s still going on because it is genuinely tough. And remember that the feeling of finally solving it is always worth the storm of frustration and patience and mild existential crisis or whatever mix of emotions it takes to get there! 
Yes winning is great but it’s not all about the winning… think aobut how many ppl go through their without doing anything, teh fact that you’re struggling with somethign you care enough aobut, and probably missing homework and not paying attention and lessons and god knows what for this challenge. Call this the clique of value the process more than the outcome, but it is in fact true. With performance comes pressure to do well, as with anything, but don’t let rip the joy of learning out of you. 
That’s a beautiful thing, it’s an amazing experince to have. Cherish that. Life would be so much more dull without challenges like these, so the fact that you’re choosing to engage in this frustrating and rewarding activity means a lot. 
How amazing is it to be part of something bigger, by being part of this journey. What to do when you’re stuck with ideas for thing. I can’t like the stress of wanting to do well does take the joy out of things… 
Also would be nice to include analogies from sport or something like that… like how Alyson Liu in skating for process over victory- 2026 Winter Olympics. Relate things to life. Like common sayings in sport, food is fuel. Don’t lose your sleep over it kinda thing? Also when working on NCC- my codebreaking hoodie thing for 10B! Helps with takign breaks.. 
Section on takign breaks, sleep and other things and time flies. Link back to codebreakers in Bletchley having hobbies. 